\section{Erasing and re-setting the text}\label{sec:erase}

For each block that is kept for translation, the old glyphs are removed from the picture and the translation is drawn
into the place they left. The stage has three parts: a mask of the glyph pixels with an estimate of the background and the
text colour, a fill (flat, or learned where the background is busy), and a fit of the translation to the room the old text had.

\subsection{Mask, background and colour}

For every line of the block the stage crops the line's rectangle with a ring of $r=\max(3,\lfloor0.18\,h\rfloor)$ pixels around it,
where $h$ is the line height. The ring is the background sample. Let $\ell(p)$ be the luma of pixel $p$ and $\tilde\ell$ the median
luma of the ring. The noise of the background is the robust spread of the ring,
\begin{equation}\label{eq:mad}
\hat\sigma\;=\;1.4826\;\median_{p\in\text{ring}}\bigl|\ell(p)-\tilde\ell\bigr| ,
\end{equation}
where the factor makes the median absolute deviation consistent with the standard deviation of a Gaussian \citep{mad}. A median
instead of a mean keeps a box edge, or a bit of another element of the page that falls in the ring, from skewing the estimate.
The background is \emph{flat} when $\hat\sigma<6$. The background colour is the per-channel median of the ring.

Glyph pixels are one of the two classes of Otsu's threshold \citep{otsu},
\begin{equation}\label{eq:otsu}
t^\star\;=\;\arg\max_t\ \omega_0(t)\,\omega_1(t)\,\bigl(\mu_0(t)-\mu_1(t)\bigr)^2 ,
\end{equation}
computed on the luma histogram of the text rectangle: the class on the side away from the background. A pixel is a candidate when
it is in that class and $|\ell(p)-\tilde\ell|\ge\max(12,\lfloor2.5\,\hat\sigma\rfloor)$. Candidates are grouped into 8-connected components, and a
component is kept only when it reaches into the text rectangle and does not run off the outer edge of the crop where a ring
exists: glyphs sit inside the ring, while box borders and rules continue past it. The mask is dilated by
$\max(1,\lfloor0.07\,h\rfloor)$ pixels to swallow the anti-aliased edges, and enclosed holes are filled. If segmentation fails (the
glyph class covers more than $0.6$ of the text rectangle, or there is none) the whole rectangle is the mask.

The text colour is the mean colour of the half of the glyph pixels that contrast most with the background, because edge pixels
are blended with it. Outlined text, such as video subtitles, has two tones, and either one alone would leave the other behind:
the mask then takes both, and the outline colour is kept for the translation. Coloured frames drawn round glyphs, such as the
box a viewer draws round a link, are taken into the mask when they are thin outlines that hold glyph pixels inside.

\subsection{Fill and inpainting}

If the background of every line of the block is flat, or inpainting is turned off, the masked pixels are filled with the
background colour. Otherwise the patch is repaired with MI-GAN \citep{migan}, a generator designed for inpainting on mobile
devices. The model takes a $4\times512\times512$ tensor whose first channel is $M-\tfrac12$, with $M=1$ for known pixels and $M=0$ for the
hole, and whose other channels are the RGB values scaled to $[-1,1]$ and multiplied by $M$. Its output is a full-frame
reconstruction $G$ that also differs from the input outside the hole, so the result is composited:
\begin{equation}\label{eq:composite}
\hat I\;=\;M\odot I\;+\;(1-M)\odot G\bigl(M-\tfrac12,\ M\odot I\bigr).
\end{equation}
A rectangle larger than a window is tiled with $512\times512$ windows at a step of $\tfrac34\cdot512=384$ pixels, each window centred on its
tile and taking real frame pixels as context; only hole pixels inside the tile are written.

Core ML runs the model on the CPU, the GPU or the Neural Engine. The inpaint bench measured the three on a $512^2$ patch
(Table~\ref{tab:migan}), and the application selects the GPU: the default setting that lets Core ML choose among all units was four
times slower there. The 14.2\,MB model is not loaded at launch; it loads at the first fill, and is released after 60\,s without one.

\begin{table}[t]
\centering
\caption{MI-GAN $512^2$ prediction time by compute unit, one patch, median of 20 timed runs after 3 warm-up runs (inpaint bench,
M4 Pro).}\label{tab:migan}
\small
\begin{tabular}{@{}lrr@{}}
\toprule
Compute unit & Median (ms) & p90 (ms)\\
\midrule
CPU only & 212.5 & 320.5\\
CPU and GPU & 13.3 & 13.4\\
CPU and Neural Engine & 20.0 & 20.3\\
All units & 55.3 & 64.1\\
\bottomrule
\end{tabular}
\end{table}

\subsection{Re-setting the translation}

The translation is laid out in the room of the old text, and the room can be larger than the old text: free background to the
right of the block is part of the patch, up to $80\%$ of the block's width and up to the edge of its column, where free means
that no pixel differs from the background colour by more than $22$.

\paragraph{Size.}
The size of the old text is read from the height $h_g$ of its glyphs: the em size is $h_g/0.88$ for Chinese, Japanese and
Korean text, whose glyphs fill about $88\%$ of the em, and for Latin text $h_g/0.72$ (cap height) or $h_g/0.95$ when the text has
descenders. The translation is then set at the largest size that fits the room. With $\mathrm{measure}(s)$ the size of the laid-out
text at font size $s$,
\begin{equation}\label{eq:fit}
s^\star\;=\;\max\bigl\{\,s\in[s_{\min},s_{\max}]\ :\ \mathrm{measure}(s)\le\mathrm{box}\,\bigr\},
\end{equation}
which is found by bisection to a resolution of $0.5$ pt (\code{TextFitter.fitFontSize}); when nothing fits, $s_{\min}$ is returned and
the caller wraps the text or lets it overflow. As a last resort a translation that does not fit at the smallest size is cut at a
word boundary and ends in an ellipsis. The ellipsis is counted by the paper bench (Section~\ref{sec:measure}), because it is a
loss: the person reads a shorter sentence than the engine returned.

\paragraph{Place and colour.}
The text starts where the old text started: the patch reaches the erase margin $r$ past the original on every side, and the text
area is the original block. The text colour is the sampled colour above, the background is the repaired patch, and Hangul syllables
of a word are joined with a word joiner so that the platform's line breaker does not split a word between two syllables.
